% ====================================================================
% Section: AI Agent Improvement Process
% ====================================================================
\section{AI Agent Improvement Process}

% --- Slide: Architecture Overview ---
\begin{frame}{Agent Architecture}
    \textbf{A two-layer knowledge system guides the LLM throughout the project.}
    \vspace{0.3cm}

    \begin{columns}[T]
        \column{0.48\textwidth}
        \textbf{Agent (Global Knowledge)}
        \begin{itemize}
            \item Source-of-truth schema
            \item Non-negotiable business rules
            \item Global consistency rules
            \item Verification behavior
        \end{itemize}

        \column{0.48\textwidth}
        \textbf{Skills (Section-Specific)}
        \begin{itemize}
            \item Methodology \& checklists
            \item Verification procedures
            \item Common mistakes
            \item Anti-hallucination rules
        \end{itemize}
    \end{columns}

    \vspace{0.4cm}
    \centering
    \small\textit{7 sections $\times$ 3 rounds each = 21 experiment iterations}
\end{frame}

% --- Slide: Iterative Workflow ---
\begin{frame}{Iterative Improvement Workflow}
    \textbf{Each section undergoes a 3-round experiment cycle:}
    \vspace{0.3cm}

    \begin{center}
        \begin{tikzpicture}[
            node distance=0.4cm and 1.0cm,
            box/.style={draw, rounded corners, minimum width=2.2cm, minimum height=0.7cm, align=center, font=\small},
            arr/.style={-{Stealth[length=2mm]}, thick}
        ]
            % Row 1: Main flow
            \node[box] (skill) {Build\\Skill};
            \node[box, right=of skill] (gen) {Generate\\Result};
            \node[box, right=of gen] (eval) {Evaluate\\Result};
            \node[box, right=of eval] (improve) {Propose\\Improvements};

            \draw[arr] (skill) -- (gen);
            \draw[arr] (gen) -- (eval);
            \draw[arr] (eval) -- (improve);

            % Feedback loop
            \draw[arr] (improve.south) -- ++(0,-0.6) -| (skill.south)
                node[pos=0.25, below, font=\scriptsize] {Update Skill \& Agent};
        \end{tikzpicture}
    \end{center}

    \vspace{0.2cm}
    \begin{itemize}
        \item \textbf{Round 1:} Baseline --- initial skill generates the first benchmark.
        \item \textbf{Round 2:} Lessons from Round 1 refine the skill before regenerating.
        \item \textbf{Round 3:} Accumulated lessons produce the final benchmark.
    \end{itemize}
\end{frame}

% --- Slide: Evaluation & Learning ---
\begin{frame}{Evaluation \& Learning}
    \begin{columns}[T]
        \column{0.48\textwidth}
        \textbf{Evaluation Dimensions}
        \begin{itemize}
            \item Requirement coverage
            \item Correctness
            \item Naming consistency
            \item Hallucination risk
            \item Documentation quality
            \item Maintainability
        \end{itemize}

        \column{0.48\textwidth}
        \textbf{Learning Mechanism}
        \begin{itemize}
            \item Each round produces an \texttt{ImproveXX} document
            \item Records issues, root causes, and proposed fixes
            \item Knowledge is \textbf{cumulative} --- never discarded
            \item Lessons are classified: Agent-level vs.\ Skill-level
        \end{itemize}
    \end{columns}

    \vspace{0.3cm}
    \centering
    \small\textit{Human approval is required at every step --- the agent proposes, never applies.}
\end{frame}

% --- Slide: Key Design Principles ---
\begin{frame}{Key Design Principles}
    \begin{enumerate}
        \item \textbf{Human-in-the-Loop:} The agent proposes changes; humans approve before any modification is applied.
        \vspace{0.2cm}
        \item \textbf{Separation of Concerns:} Dedicated workers for skill building, result generation, evaluation, and improvement --- no worker exceeds its role.
        \vspace{0.2cm}
        \item \textbf{Cumulative Knowledge:} Skills and improvement documents grow over rounds; previous lessons are preserved and refined.
        \vspace{0.2cm}
        \item \textbf{Anti-Hallucination:} Workers operate only on explicitly injected files; no repository exploration or hidden context is allowed.
    \end{enumerate}
\end{frame}
